\unit{Noise, loss and quantum channels}
\unitref{W5L1; Week 5 tutorial}
\subsection{Why a subsystem does not evolve unitarily}
For an isolated quantum system,
\(\rho'=U\rho U^\dagger\).
An optical mode in the laboratory interacts
with unobserved environmental degrees of
freedom: lost photons, thermally excited
modes, unwanted phase fluctuations and detectors.
If the environment starts in \(\rho_E\),
\[
\boxed{\rho'_S=
\operatorname{Tr}_E[
U_{SE}(\rho_S\otimes\rho_E)U_{SE}^\dagger].}
\]
The joint evolution is unitary but the reduced
system evolution generally is not. Tracing out
the environment is an \emph{average over
unobserved information}, not a measurement
that selects a particular outcome.

\subsection{Loss as a beam splitter with vacuum}
Introduce an initially empty environment mode \(e\):
\[
a_{\rm out}=\sqrt\eta\,a_{\rm in}
+\sqrt{1-\eta}\,e_{\rm in}.
\]
This is the physically correct quantum loss
model; replacing \(a\) by \(\sqrt\eta a\)
alone would give
\([a_{\rm out},a_{\rm out}^\dagger]=\eta\ne1\).
The environment operator supplies the vacuum
noise that preserves the commutator.
For a single photon,
\[
\ket{1,0}_{S,E}\longrightarrow
\sqrt\eta\ket{1,0}
+\sqrt{1-\eta}\ket{0,1}.
\]
Tracing over the environment gives
\[
\boxed{\rho'_S=\eta\ket1\bra1+
(1-\eta)\ket0\bra0.}
\]
The probability of having a photon is \(\eta\);
the environment is responsible for the mixture.

\subsection{Kraus operators: a system-only description}
The same channel can be written
\[
\boxed{\mathcal E(\rho)=\sum_kK_k\rho K_k^\dagger,
\quad\sum_kK_k^\dagger K_k=I.}
\]
For a single-rail state supported on
\(\{\ket0,\ket1\}\), loss has
\[
K_0=\begin{pmatrix}1&0\\0&\sqrt\eta\end{pmatrix},
\qquad
K_1=\begin{pmatrix}0&\sqrt{1-\eta}\\0&0\end{pmatrix}.
\]
The \(K_0\) branch contains no environment
excitation, while \(K_1\) represents losing
the photon. These are \emph{not} two
alternative deterministic maps; they form
one trace-preserving channel when summed.

\begin{worked}{A superposition passes through loss}
For \(\ket\psi=\alpha\ket0+\beta\ket1\),
\[
\rho'=\sum_kK_k\ket\psi\bra\psi K_k^\dagger
=\boxed{\begin{pmatrix}
|\alpha|^2+(1-\eta)|\beta|^2&
\sqrt\eta\,\alpha\beta^*\\
\sqrt\eta\,\alpha^*\beta&
\eta|\beta|^2
\end{pmatrix}}.
\]
Check the trace: it is 1.
The excited population is suppressed by
\(\eta\), but vacuum--one-photon
coherence by \(\sqrt\eta\).
For \(\eta=1\) the original pure state is
recovered. For \(\eta=0\) the output is vacuum.
\end{worked}

\subsection{Phase damping is different from photon loss}
Phase noise can suppress coherences without
changing photon-number probabilities.
For a number basis, a random phase
\(e^{i\phi\hat n}\) has no effect on
diagonal populations but averages away
off-diagonal terms such as \(\ket0\bra1\).
A simple dephasing channel acts on a qubit as
\[
\rho_{01}\mapsto\lambda\rho_{01},\quad
\rho_{00}\mapsto\rho_{00},\quad
\rho_{11}\mapsto\rho_{11},
\qquad0\le\lambda\le1.
\]
Hence direct photon counting can miss phase
noise that is obvious in homodyne detection.

\subsection{When measurement becomes conditioning}
If an environmental or ancilla outcome \(k\)
is actually observed and kept, use
\[
\tilde\rho_k=K_k\rho K_k^\dagger,\quad
p_k=\operatorname{Tr}\tilde\rho_k,\quad
\boxed{\rho_k=\tilde\rho_k/p_k}.
\]
This is a trace-decreasing \emph{conditional}
operation. Summing over all outcomes instead
recovers the unconditional channel.
For conditional photon subtraction, a weak
beam splitter followed by a one-photon
detector click implements an operator
approximately proportional to \(a\).
It weights higher Fock components by \(\sqrt n\),
not simply removing one photon deterministically.

\begin{checkpoint}
Ask which information is available.
If the environment is ignored, trace it
out or sum Kraus branches. If a detector
outcome is heralded, \emph{select and
renormalise} that branch, retaining its
separate success probability.
\end{checkpoint}
